\chapter{The Least Principle and the Unseen}
\label{chap:least-principle-and-the-unseen}

This chapter closes Part I, and by its end the instrument is fully assembled. It governs motion when nothing
visibly forces it, and it confronts the question every finite instrument must answer honestly: what can a
record say about what it did not directly see? The chapter runs on one division. The criterion is ours --- the
cost we score a motion by, the rule we ask it to make least --- but the path that actually answers, the one
configuration that makes the cost stationary, is the world's; we set the contest, the world names the winner.
The two halves --- the least principle, that the actual motion extremizes an action, and the unseen, the
stretches a record must reconstruct, the boundary that makes a solution unique, the horizon past which
prediction fails --- turn out to be one discipline under that division. The instrument fills the unseen with
its own least-cost guess, ours to compute, and the least principle is, in the end, the demand that the unseen
leave the smoothest trace it can; but what no guess of ours can fill --- the boundary that fixes the answer,
the horizon past which prediction fails --- the world keeps for itself. By the chapter's close the instrument
has every capability the rest of the book will use.

\section{The Minimizing Variations Effect: the instrument's economy}
\label{se:minimizing-variations}

Of all the paths a body could record between two fixed configurations, the instrument lays down the one that
costs the least. This is the least principle, and to the instrument it is a statement about economy: the
actual motion is the path that demands the fewest reconciliations --- the least re-filing --- to lay down
from one endpoint to the other. Writing the cost of a path as an action, a tally accumulated along it,
\begin{equation}
  J[x] = \int_a^b L\bigl(t, x, \dot{x}\bigr)\, dt,
\end{equation}
the principle says the recorded path is the one at which this cost is stationary,
\begin{equation}
  \delta J = 0 .
\end{equation}
Read this as a description, not a calculation. The cost is ours --- we wrote the rule $L$ that scores each
path --- but where it bottoms out is not. It says: nudge the path a little, holding its endpoints fixed, and to
first order the cost does not change --- the instrument is sitting at the bottom of the cost we set, a place
the world picks out and the rule only names, where small re-routings buy nothing. The condition that a path sit there everywhere along its length has a shape,
and that shape is the Euler--Lagrange equation,
\begin{equation}
  \frac{d}{dt}\frac{\partial L}{\partial \dot{x}} = \frac{\partial L}{\partial x} .
\end{equation}
This equation is not something the instrument solves to find the motion; it is the \emph{form} a least-cost
path must have, the local signature of a global economy --- ours to write down, the world's to satisfy or not.
Where it holds, the path is stationary; where it fails, a cheaper re-routing exists and the path is not the one
the instrument records.

Two consequences of this economy organize the whole book, and both belong here. The first is Noether's
theorem, met in the last chapter for rotation and stated now as the general rule: to every continuous symmetry
of the action there corresponds a conserved current. Where the cost is unchanged under sliding the clock
forward --- time-translation symmetry --- the conserved quantity is the energy; where it is unchanged under
sliding space --- space-translation symmetry --- the conserved quantity is the momentum; where it is unchanged
under turning --- rotational symmetry --- the conserved quantity is the angular momentum. The conservation
laws scattered through the rest of the book are not separate facts but one theorem, read off the symmetries
the instrument cannot tell apart: a sameness the action does not register becomes a count the motion cannot
change.

The second consequence completes the conjugate-pair structure the force chapter introduced. The Lagrangian is
a function of position and velocity, $(q, \dot q)$; the instrument's update rule, Hamilton's, runs on position
and its conjugate momentum, $(q, p)$. The bridge between the two descriptions is the Legendre transform,
\begin{equation}
  H(q, p) = p\,\dot q - L(q, \dot q),
\end{equation}
which trades the velocity $\dot q$ for the conjugate momentum $p = \partial L/\partial \dot q$ and returns the
Hamiltonian. The two formalisms are one economy seen from two sides --- the Lagrangian extremizing a cost
along whole paths, the Hamiltonian stepping the conjugate pair forward point by point --- and the Legendre
transform is the change of variables that carries each into the other. With it the analytical mechanics begun
in the force chapter is complete: a least principle, a conserved count for every symmetry, and a conjugate
pair the next parts will quantize.

The economy is visible on the bench. Offer the instrument two candidate records between the same endpoints
--- the true motion and a slight detour, a path that wanders a little wider and then hurries to make up the
time. Tally the cost along each: the detour, being longer and more sharply turned, demands more
reconciliations to lay down, and the instrument, asked which record it would actually produce, returns the
cheaper one. We supplied both candidates and the cost that ranks them; which one comes back is the world's
verdict, not ours. The least principle is not a mystical preference for economy imposed on nature from
outside; it is the plain fact that the record the instrument lays down is the one that took the fewest
distinctions to write.

The honest floor is a smooth-shadow analogy. What the instrument actually secures is a discrete least-cost
path over a sampled record --- the cheapest sequence of counts between the endpoints --- and the
Euler--Lagrange equation is the smooth shadow that discrete economy casts as the sampling is refined. The
bridge between the cheapest sequence of counts and the differential equation is an analogy, ours to draw and
named as one. The reading is that the law of motion is not a force imposed from outside but an economy obeyed
from within: the rule of frugality is ours to write, and the path that obeys it --- the one the instrument
actually lays down --- is the world's.

\section{The Repeatability of Invisible Motion Effect: the smoothest honest guess}
\label{se:repeatability}

The least principle has an immediate use: it tells the instrument how to fill a stretch it did not see.
Suppose a record shows two configurations and, between them, no events at all. To a finite-resolution
instrument this silence is not ignorance but information --- any acceleration or inflection large enough to
register would have left a trace, so whatever happened between the anchors sat below the floor. The
instrument must reconstruct the unseen stretch. What the world did down there is the world's, and the
instrument will not claim to have seen it; but it must lay down \emph{some} continuation, and it does so with
its smoothest honest guess: the minimal-curvature completion, the curve that bends as little as the endpoints
allow. That curve is the one whose fourth derivative vanishes,
\begin{equation}
  \Psi^{(4)} = 0,
\end{equation}
with neighboring patches glued so that the curve, its slope, and its bend are all continuous --- a $C^2$ fit,
the natural cubic spline through the anchors. The guess is ours; that it is forced to be \emph{this} guess and
no other is the discipline's doing. And this smoothest guess is exactly the least-action extremal of the
previous section: smooth dynamics are precisely the histories that leave no trace, so two observers who see
nothing between the same endpoints reconstruct the very same cubic.

The honest floor is again a smooth-shadow analogy --- the discrete minimal-curvature reconstruction is the
shadow of the continuum extremal, ours to draw and named as one. The reading is that invisible motion is not
unknown but, as a reconstruction, uniquely determined: the \emph{guess} is fixed even where the world's doing
is not. The unseen stretch is filled by the smoothest path the endpoints allow, which is the least-action
path, which is exactly why a motion that leaves no trace is repeatable at all: every honest instrument, asked
to fill the same silence, fills it the same way --- not because the world told them what happened, but because
the discipline hands them all the same least-committed guess. The instrument does not guess wildly into the
dark; it lays down the one curve its economy demands, and labels it for what it is --- a reconstruction, the
smoothest the anchors permit, not a thing seen.

\section{The Dirichlet--Bancroft Effect: an interior fixed from its edge}
\label{se:dirichlet-bancroft}

Here the boundary reconciliation the momentum chapter introduced returns as a power: the instrument's ability
to fix a whole interior from its edge alone. A model can be exact and internally consistent and still fail to
pin down a unique answer, because the governing equations by themselves leave room for more than one; what
selects the true answer is boundary data committed to the record. The standing bench instance is the global
positioning system, and Bancroft's method of reading a position from it. The satellite positions and timing
are fully specified, and yet the equations that locate a receiver are quadratic --- they admit two solutions,
a true position and a phantom. The interior is under-determined; the algebra alone cannot say which is real.

Picture it as trilateration. Each satellite tells the receiver only its distance --- a sphere of possible
positions centered on that satellite. Two such spheres meet in a circle, three in a pair of points: the
receiver is at one of them, but the ranges alone do not say which. The interior of the problem, the
receiver's location, is pinned to two candidates and no further. It is the edge that breaks the tie --- the
constraint that the receiver sits on or near the Earth's surface, a condition imposed from the boundary,
selects the one point of the pair where a receiver can actually be.

A single boundary fact decides it. The constraint that the receiver lies near the Earth's surface --- a fact
about the edge, not the interior --- discards the phantom and collapses the two solutions to one. The boundary
data is ours to bring: which satellites we range to, which surface constraint we commit to the record, is the
hand at the rim. What follows from it is not ours --- the interior, where the receiver actually is, the world
fixes uniquely once the edge is set. The instrument has fixed that interior from the boundary alone, without
ever needing to scan the interior directly: it stands at the edge, reconciles what the edge demands, and the
interior follows. This is the same boundary reconciliation that will, in the quantum and relativity parts, fix
a record's order from its frontier without opening the record up; here it appears in its plainest classical
form, as trilateration.

The trilateration is the plainest instance, but the structure is general, and naming it sharpens the point.
When the interior of a region carries a field with no sources of its own --- a potential, a steady
temperature, an equilibrium configuration --- that field is harmonic: it satisfies Laplace's equation,
\begin{equation}
  \nabla^2 \Psi = 0,
\end{equation}
and a harmonic field is fixed completely and uniquely by its values on the boundary. The interior holds no
freedom the edge has not already settled. The reason is the maximum principle: a harmonic field attains its
largest and smallest values only on the boundary, never strictly inside --- there is no interior peak or
valley the boundary did not put there. So to know the field everywhere within, the instrument needs only its
values on the frontier --- those values ours to gather at the rim --- and the bulk is the boundary's unique
harmonic continuation inward, the world's to fill in. This is the precise form of fixing an interior from its
edge, the same structure --- a frontier that determines a bulk --- that the boundary reconciliations of the
later parts carry.

The honest floor is a finite count obligation: the count is the number of admissible solutions, and the
boundary is the obligation that reduces it from two to one. The reading is that completeness is not
uniqueness. What a record must commit to is its boundary --- ours to supply --- and with it the interior is the
world's to determine; without it the equations, however exact, leave a fork the instrument cannot close, and
no rule of ours closes it.

\section{The Butterfly Effect: the instrument's horizon}
\label{se:butterfly}

The last of Part I's effects gives the instrument its own horizon. Prediction has a limit, and the usual name
for it is chaos --- sensitive dependence, a butterfly's wing swaying a storm. But the instrument names the
limit differently, and more honestly: it is not, in the first place, the dynamics amplifying a difference
exponentially, but the instrument's own floor, grown until it swamps the signal. The wish is ours --- to
forecast as far ahead as we like --- and the floor is ours too; what the world contributes is the growth that
turns the floor into a wall. Two histories that differ only below the resolution floor are, to the instrument,
one history; as time runs forward, the gap between such histories grows, and at some finite depth it exceeds
what the record can hold. Count the admissible microhistories consistent with a
single coarse reading: that count climbs as finer and finer distinctions go unrecorded, and the horizon of
prediction is the moment it overruns the instrument's capacity to carry it. Beyond that horizon the instrument
cannot predict --- not because the future is unknowable in principle, but because the record cannot store the
sub-resolution information that would decide it.

The amplification has a rate, and naming it shows how unforgiving the horizon is. Even granting the dynamics
its divergence rate --- which is the world's to set --- what the instrument secures is its own horizon: a gap
that starts at the floor, of size $\epsilon_0$, ours, and grows, where the motion is chaotic, exponentially,
\begin{equation}
  \delta(t) \sim \epsilon_0\, e^{\lambda t},
\end{equation}
with $\lambda$ the Lyapunov exponent, the rate at which nearby histories diverge. The amplified floor reaches
the scale at which the record can no longer hold it after a time
\begin{equation}
  t_h \sim \frac{1}{\lambda}\,\ln\frac{1}{\epsilon_0},
\end{equation}
the prediction horizon. The cruelty is in the logarithm: because the horizon grows only as the logarithm of
the precision, buying a thousandfold finer floor buys only a few more multiples of $1/\lambda$ before the
floor swamps the signal again. The instrument cannot push its horizon out by force; the divergence converts
each decimal of added precision into a fixed, small extension of foresight, and no finite floor reaches past
the wall.

The honest floor is a finite count obligation: the microhistory count measured against capacity, the depth at
which it overflows. The reading deserves to be stated flatly, because it is easy to mistake. The butterfly's
horizon belongs to the apparatus, not the world. It is a combinatorial limit on what a finite record can
distinguish, not an exponential sensitivity of the equations; one is a fact about an instrument, the other a
claim about the dynamics, and only the first is what the experiment secures. The instrument knows exactly how
far it can see; the forecast it would like to give without end, it gives only this far, and it does not pretend
to see past its own horizon.

\bigskip

With that, Part I closes, and the instrument is complete. Across four chapters it has been assembled piece by
piece: it records distinctions as a count; it has a resolution floor below which it reads nothing; it
reconciles records at their boundaries, fixing an interior from its edge; it carries a tally around a closed
loop and reads the holonomy where the loop returns; it lays down the path of least cost, spending the fewest
distinctions; and it knows its own horizon, the depth past which its floor swamps the signal. Six
capabilities, and four honest floors --- a finite count, a smooth shadow, a no-go wall, and, in name if not
yet in force, a label --- with which it has read all of mechanics off a finite record, drawing in every effect
the one line between what the world fixed and what was the instrument's own to lay down. Everything that
follows is this same instrument, unchanged, carried to larger and larger scales: to the fields, where its
holonomy becomes a residue it proves; to heat, where its count becomes a temperature; to the quantum, where
its count becomes an amplitude; to relativity, where its floor and its loop bend space and time; and to the
foundations, where it is turned at last upon itself. The instrument is built. Part II sets it among the
fields.

\bigskip
\begin{center}
\rule{0.4\textwidth}{0.4pt}
\end{center}
\bigskip

\noindent Offered many accounts of one motion, all fitting the facts already entered, the record the
instrument lays down is the shortest admissible one --- the account that spends the fewest distinctions to
reach from the first configuration to the last. The standard of economy is ours: we set the cost each account
is scored by, the measure of what counts as spare. Which account is genuinely cheapest --- which one, of all we
propose, the world realizes --- is not ours to decree. A case is built on that cheapest account and no other:
the most economical reading that fits every count already filed --- not the fullest telling, nor the most
cautious, but the leanest account that still answers for every mark entered.

The wheel does not stop being counted where the sensor stops looking. Between two passes the instrument did not
catch, the wheel still turned; and the record must say what it did there without pretending to have watched it.
So the instrument lays down the least-committed going-around the anchoring counts allow --- the smoothest
continuation of the turning, bending as little as the two ends permit --- and marks it for exactly what it is: a
reconstruction, the frugal guess the discipline forces, never a pass it witnessed. The turning between the
anchors is the world's; the curve drawn through the gap is ours, and it is labeled a guess the moment it is
entered. The instrument reports what it saw as seen and what it filled as filled, and it does not let the two
change coats. And the guess is no single instrument's private invention: any honest instrument asked to fill
the same silence fills it the same way --- not because the world reported what happened down there, but because
the discipline hands each of them the very same least-committed continuation.

With that the car's own instrument is complete. Across Part I it has been built piece by piece: it counts; it
has a floor below which two states read as one; it reconciles an interior from its edge, closing from the
boundary a fork the interior alone leaves open; it carries a tally around a closed loop and reads the residue
where the loop comes home; it lays down the least-cost account, the one that spends the fewest distinctions;
and it knows its own horizon, the depth past which its own floor, grown, swamps the reading --- a limit that
belongs to the apparatus and not to the world, on what a finite record can tell apart, never a failure of the
world to be definite. Six capabilities --- and, standing under them, four honest floors it will never
overclaim: a finite count it can check, a smooth shadow it draws through the gaps but never holds, a wall it
records and cannot cross, and --- in name if not yet in force --- a label, the fourth floor named here and not
yet made to bear weight, reserved for a reckoning much further on. The hall-speedometer is fully assembled;
nothing after this adds a part to it, only carries the same instrument into a larger field.

The charge it enters against the record is scored, now, by that same economy. Of all the accounts that fit the
counts, the charge the instrument stands behind is the cheapest admissible one --- the loop count read the most
frugal honest way, not one distinction more than the record compels. And the trace, what the record retains of
the turning, is now exactly that minimal retained account: not every wandering the wheel might have made, but
the least its anchors can be made to carry. So the payment, when it is finally named, will be the least the
record can honestly demand --- no more than the cheapest account will bear, and no less. The chapter need not
set the figure; it is enough that what will be owed is fixed to the frugal reading, an economy before it is
ever a sum. What the record will make anyone answer for is not padded by a fear of missing something, nor
discounted by a wish to owe less; it is the one figure the counts themselves compel, read as sparely as
honesty allows.

And the demand on the three instruments sharpens here in a way it has not before. Each of the three fills its
own unseen with its own least-cost guess; each reconstructs the stretches it did not watch. The record is
consistent under the law only if the three land on the same number even there --- where every one of them has
inferred rather than seen, the cheapest accounts must still agree. It is one thing for three instruments to
agree on what all three witnessed; it is a stronger thing for them to agree on what each only reconstructed in
the dark. That stronger agreement is what the case will in the end demand of them.

The instrument is built --- but the record was never the instrument's alone, and Part I has used only the first
of the three hands the case needs. The second, named since the opening, does not ride the wheel at all. It
stands off from the car and reads it from outside, by a field that reaches the car where the wheel cannot ---
a reach that touches the moving body without any part of the instrument riding it. Part I read the car from
within, off the turning of its own wheel; the next part reads it from without, off a field that never touches
the gear train. Two readings of one speed,
then --- one taken from inside the motion and one from outside it, kept in hands that share no single part ---
and the case will hold only if the two are made to meet on one number. That is the instrument the coming
chapters bring to the bench, and the wave is the first field it meets. The account has been made as
cheaply as it honestly can; the record that must close is still open, and a second hand now comes to write in
it. The court does not sit until the record is complete.
